%\documentclass[a4paper,11pt]{article}
\documentclass[11pt]{amsart}

\usepackage[dvipsnames]{xcolor}
\usepackage{todonotes}
\usepackage[utf8]{inputenc}
\usepackage{comment,enumerate,graphicx}
\usepackage{enumitem,amsmath,amssymb}
\usepackage{amsthm}
\usepackage{thmtools}
\usepackage{url}
\usepackage{caption}
\usepackage{subcaption}
\usepackage[top=23 mm, bottom=23 mm, left=21 mm, right= 21 mm]{geometry}
\usepackage{relsize}
\usepackage{microtype}


\usepackage[hidelinks]{hyperref}
%%%% Code for property numbering in theorems.
\makeatletter
\newcommand{\labelinthm}[1]{%
   \label{temp#1}
   \protected@write \@auxout {}{\string \newlabel{#1}{{\emph{\ref{temp#1}}}{\thepage}{\emph{\ref{temp#1}}}{temp#1}{}} }%
   %\hypertarget{temp#1}{\emph{\ref{temp#1}}}
}
\makeatother
\usepackage{adjustbox}
\usepackage{tikz}
\usetikzlibrary{math,calc,intersections, scopes}
\usetikzlibrary{positioning,arrows,shapes,decorations.markings,decorations.pathreplacing, decorations.pathmorphing,matrix,patterns}
\tikzstyle{vertex}=[circle,draw=black,fill=black,inner sep=0,minimum size=5pt,text=white,font=\footnotesize]
\tikzstyle{redvertex}=[circle,draw=red,fill=red,inner sep=0,minimum size=5pt,text=white,font=\footnotesize]
\definecolor{amber}{rgb}{1.0, 0.75, 0.0}
\definecolor{darkgreen}{rgb}{0.18, 0.7, 0.46}

\newcommand\anote[1]{\todo[color=Green!50]{#1}}
\global\long\def\am#1{\textcolor{red}{\textbf{[AM comments:} #1\textbf{]}}}
\global\long\def\zh#1{\textcolor{blue}{\textbf{[ZH comments:} #1\textbf{]}}}

\declaretheorem[name=Theorem]{theorem}


\newtheorem{lemma}[theorem]{\bf Lemma}
\newtheorem{corollary}[theorem]{\bf Corollary}
\newtheorem{claim}[theorem]{\bf Claim}
\newtheorem{proposition}[theorem]{\bf Proposition}
\newtheorem{conjecture}[theorem]{\bf Conjecture}
\newtheorem{question}[theorem]{\bf Question}
\newtheorem{problem}[theorem]{\bf Problem}
\newtheorem*{theorem*}{\bf Theorem}

\theoremstyle{definition}

\newtheorem{remark}[theorem]{\bf Remark}
\newtheorem{definition}[theorem]{\bf Definition}
\newtheorem{result}[theorem]{\bf Result}
\newtheorem{example}{\bf Example}
\newtheorem{exampletwo}{\bf Example}
\renewcommand\theexample{\textbf{E\arabic{example}}}
\renewcommand\theexampletwo{\textbf{F\arabic{exampletwo}}}
\newcommand\eref[1]{\emph{\ref{#1}}}

\newenvironment{claimproof}{%
  \renewcommand{\qedsymbol}{$\boxdot$}%
  \begin{proof}%
}{%
  \end{proof}%
  \renewcommand{\qedsymbol}{$\square$}%
}


\newcommand\claimproofend{\renewcommand{\qedsymbol}{$\boxdot$}
\end{proof}
\renewcommand{\qedsymbol}{$\square$}
}
\def\eps{\varepsilon}
\def\vx{\textbf{x}}
\def\vy{\textbf{y}}
\def\vz{\textbf{z}}
\def\Hom{\mathrm{Hom}}
\def\cC{\mathcal{C}}
\def\cD{\mathcal{D}}
\def\cF{\mathcal{F}}
\def\cG{\mathcal{G}}
\def\cH{\mathcal{H}}
\def\cN{\mathcal{N}}
\def\cP{\mathcal{P}}
\def\cQ{\mathcal{Q}}
\def\cS{\mathcal{S}}
\def\bF{\mathbb{F}}
\def\bE{\mathbb{E}}
\def\bR{\mathbb{R}}
\def\bZ{\mathbb{Z}}
\def\Pb{\mathbb{P}}
\def\ex{\mathrm{ex}}

\def\be{\boldsymbol{e}}
\def\bx{\boldsymbol{x}}
\def\by{\boldsymbol{y}}
\def\bz{\boldsymbol{z}}

\DeclareMathOperator{\red}{red}
\DeclareMathOperator{\blue}{blue}
\DeclareMathOperator{\spn}{span}

\title{\vspace{-0.9cm} On the Probability a Weighted Bernoulli Sum Exceeds Its Mean}
\date{}

\author{Aleksa Milojevi\'c}
\author{Benny Sudakov}
\thanks{Department of Mathematics, ETH Z\"urich, Switzerland. Email: {\tt \{aleksa.milojevic, benjamin.sudakov\}@math.ethz.ch}. Research supported in part by SNSF grant 200021-228014.}

\begin{document}

\maketitle

\begin{abstract}
Let $w_1, \dots, w_m$ be positive real weights summing to $1$, and let $v_1, \dots, v_m$ be i.i.d. Bernoulli$(p)$ random variables. If we let $X=\sum_{i=1}^m w_i v_i$, then we ask for which $p\in (0, 1)$ we have that
\[\Pb\big[X\geq \bE[X]\big]\geq p?\] 
In this short note, we observe a connection of this question with a version of the Manickam--Mikl\'os--Singhi conjecture, which allows one to resolve it for sufficiently small values of $p$.
\end{abstract}

\section{Introduction}

Let $w\in \bR_{>0}^m$ be a vector of positive weights satisfying $\sum_{i=1}^m w_i=1$, and let $v\in \{0, 1\}^m$ be a random binary vector, in which every coordinate equals $1$ independently and with some fixed probability $p$. Then, the random variable $X=\langle w, v\rangle$ takes
values in $[0,1]$ and satisfies $\bE[X]=p$. 

In this note, we study lower bounds on the probability that $X$ exceeds $p$, i.e. its mean. In general, one cannot expect a lower bound better than $p$, because of the following example. If $m=1$ and $w_1=1$, then $X=v_1$ and $\Pb[X\ge \bE[X]]=\Pb[v_1=1]=p$. The following question asks for which value of $p$ this example is indeed the worst one. 

\begin{question}\label{q:prob}
For which values of $p\in (0, 1)$ does the following hold? For every $w\in \bR_{>0}^m$ such that $\sum_{i=1}^m w_i=1$, if $v\in \{0, 1\}^m$ is a random vector with i.i.d. Bernoulli$(p)$ coordinates, then
\[\Pb[\langle v, w\rangle\geq p]\geq p.\]
\end{question}

To the best of our knowledge, the first result proved in this direction comes from the work of Alon, Emek, Feldman and Tennenholtz~\cite{AEFT13} on information leakage in adversarial games, where they have shown that $\Pb[X\ge \bE[X]]\ge p(1-p)/10$ for all $p$, which is tight up to a constant factor for $p\leq 1/2$. One natural guess would be that Question~\ref{q:prob} holds for all $p\leq 1/3$ and for $p=1/2$. In Proposition~\ref{prop:counterexample} we show that the inequality does not hold for other values of $p$.

It is also worth noting that \cite{AEFT13} already proposes a very elegant coupling argument showing that the answer to Question~\ref{q:prob} is yes if $p$ is of the form $p=\frac{1}{n}$ for some positive integer $n$. In fact, the proof in this note is inspired by their argument. %Also, the reason for the $p\leq 1/3$ assumption is that when $p>1/3$ the statement is simply not true (unless $p=1/2$), as we will show in Proposition~\ref{prop:counterexample}.

Let us also note that in the special case where all weights are equal (so that $mX\sim \mathrm{Bin}(m,p)$), related lower bounds for $\Pb[\mathrm{Bin}(m,p)\ge mp]$ have been studied; see, e.g., Greenberg--Mohri~\cite{GM14}, Pelekis--Ramon~\cite{PR16}, and Doerr~\cite{Doe18}. Finally, let us mention that Question~\ref{q:prob} is also vaguely reminiscent of Feige's conjecture, which states that for independent positive random variables $X_1, \dots, X_n$, each of which has expectation $1$, we have $\Pb[X_1+\dots+X_n\geq n+1]\leq 1-1/e$. 

The purpose of this note is to explain an unexpected (at least to us) connection of this problem to the classical
Manickam--Mikl\'os--Singhi (MMS) conjecture on the number of nonnegative $k$-sums.

\begin{conjecture}\label{conj:MMS}
Let $n,k$ be positive integers, with $n\ge 4k$. For any real numbers $x_1,\dots,x_n$ whose sum is $0$, at least $\binom{n-1}{k-1}$ of the $k$-element subsets $S\subseteq [n]$ satisfy
\[
  \sum_{i\in S} x_i \ge 0.
\]
\end{conjecture}

This question was first raised by Bier and Manickam~\cite{BM87}, who proved Conjecture~\ref{conj:MMS} when $k$ divides $n$ or when $n$ is much larger than $k$, and this exact form of the conjecture was later put forward by Manickam and Mikl\'os~\cite{MM88} and also by Manickam and Singhi~\cite{MS88}. In 2012, Tyomkyn~\cite{Tyom12} significantly improved the quantitative bounds, proving the conjecture for $n>k(4e\log k)^k$. Later, Alon, Huang and Sudakov~\cite{AHS11} proved the polynomial bound and showed the conjecture when $n\ge 33k^2$. Finally, a linear bound $n\geq 10^{46} k$ was obtained by Pokrovskiy~\cite{Pok15}.

% A condition of the form $n\geq 3k+2$ is necessary for Conjecture~\ref{conj:MMS} to hold, as can be seen from the following example of Bier and Manickam~\cite{BM87}. If $n=3k+1$ and $x_{1}=\cdots=x_{3k-2}=3$ and $x_{3k-1}=x_{3k}=x_{3k+1}=2-3k$, then the only $k$-tuples with nonnegative sum are those which include only positive numbers, and therefore there are at most $\binom{n-3}{k}< \binom{n-1}{k-1}$ such $k$-tuples. 

% Our goal is to show Conjecture~\ref{conj:prob} if the following variant of the Conjecture~\ref{conj:MMS} holds.

% \begin{question}\label{q:MMS_variant}
% Is the following statement true?

% For any constants $\eps, C>0$, and every sufficiently large integers $n, k$ satisfying $n\geq (3+\eps)k$, we have the following. Let $x_1, \dots, x_n$ be real numbers whose sum is $0$, such that at most $C$ among them are nonnegative and such that any two negative values $x_i, x_j$ are equal. Then, at least $\binom{n-1}{k-1}$ of the $k$-element subsets $S\subseteq [n]$ have nonnegative sum.
% \end{question}

The main purpose of this note is to prove the following theorem. This theorem, combined with the result of Pokrovskiy~\cite{Pok15}, shows that the answer to Question~\ref{q:prob} is positive for all $0\leq p< 10^{-46}$. Moreover, it shows that conditionally on the Manickam-Mikl\'os-Singhi conjecture being true, Question~\ref{q:prob} holds for all $0\leq p\leq 1/4$.

\begin{theorem}\label{thm:main}
If there is a constant $c$ such that Conjecture~\ref{conj:MMS} holds for all $n\geq ck$, then Question~\ref{q:prob} holds for all $p\in [0, 1/c)$.
\end{theorem}

Interestingly, the same theorem would still hold true if one weakens Conjecture~\ref{conj:MMS} to only ask about sequences $x_1, \dots, x_n$ where all negative entries have the same value.

\section{Proofs}

\begin{proof}[Proof of Theorem~\ref{thm:main}.]
We split this proof into two parts - first, we will show the positive answer Question~\ref{q:prob} for all rational numbers $p\in [0, 1/c)$, and then we will extend it to all $p\in [0, 1/c)$ via a simple limiting argument.

Thus, let us assume that we have $p=\frac{k_0}{n_0}$, a rational number, and let $w\in \bR_{>0}^m$ be the corresponding weight vector. Let $t$ be a large integer, and let us set $k=k_0t, n=n_0t$.

The crucial object of the proof is the random $0/1$-matrix $M$, with $n$ rows and $m$ columns. To generate $M$, we will fill out every column with $n-1$ zeros and exactly $1$ one, with the position of the $1$ being uniformly picked among the $n$ possible rows, independently of all other decisions. If we denote the rows of $M$ by $r_1, \dots, r_n$, we calculate the inner products $y_i=\langle w, r_i\rangle$, and we choose a $k$-element set $S\subset [n]$ uniformly at random. Finally, we set $Y=\sum_{i\in S}y_i$. 

\begin{claim}
The random variable $Y$ has the same distribution as the inner product $X=\langle w, v\rangle$.
\end{claim}
\begin{claimproof}
Since $Y=\langle w, \sum_{i\in S}r_i\rangle$, it suffices to show that $\sum_{i\in S}r_i$ follows the same distribution as $v$. For any fixed set $S$ and $j\in [m]$, the $j$-th coordinate of $\sum_{i\in S}r_i$ is $1$ precisely if one of the rows corresponding to $S$ contains a $1$ in the $j$-th column, which happens exactly with probability $\frac{|S|}{n}=\frac{k}{n}=p$. Since the distribution of $\sum_{i\in S}r_i$ is the same as the distribution of $v$ for any fixed set $S$, the same is true if $S$ is chosen uniformly at random.
\end{claimproof}

\begin{claim}
We have \[\Pb[Y\geq p]\geq p.\]
\end{claim}
\begin{claimproof}
Let us define $z_i=y_i-\frac{1}{n}$. First, observe that $y_i$ either takes value $0$ (if all coordinates of $r_i$ are $0$), or it takes value at least $\min_{j\in [m]}\{w_j\}>0$. Further, observe that if $t$ is sufficiently large, then $\frac{1}{n}=\frac{1}{n_0t}<\min_{j\in [m]}\{w_j\}$. Therefore, the only distinct negative value which appears among $z_i$'s is $-\frac{1}{n}$.

Also, since the sum of coordinates of $w$ is $1$ and $M$ contains exactly one $1$ per column, we have \[\sum_{i=1}^n y_i=\langle w, \sum_{i=1}^n r_i\rangle=\langle w, \mathbf{1}\rangle=1.\]
Thus, $\sum_{i=1}^n z_i=0$. 

Finally, since $p=\frac{k}{n}<1/c$, we have $n\geq ck$. Thus, if Conjecture~\ref{conj:MMS} holds for $n\geq ck$, we can apply it to the numbers $z_1, \dots, z_n$ to conclude that a uniform random set $S\subset[n]$ of size $k$ has nonnegative sum with probability at least $\binom{n-1}{k-1}/\binom{n}{k}=\frac{k}{n}=p$. 

Since $\sum_{i\in S}z_i\geq 0$ is equivalent to $Y=\sum_{i\in S} y_i\geq \frac{|S|}{n}=p$, we conclude that with probability at least $p$ over the choice of $S$ we have $Y\geq p$, as we wanted to show.
\end{claimproof}

Combining the two claims above shows that $\Pb[\langle v, w\rangle\geq p]=\Pb[Y\geq p]\geq p$ whenever $p$ is a rational and at most $\frac{1}{c}$. To show the conjecture when $p$ is irrational, let $p_i$ be a sequence of rationals approaching $p$ from below and observe that $\Pb[\langle w, v\rangle \geq p_i]\geq p_i$ (where the entries of $v$ are still distributed according to Bernoulli$(p)$, as before). This inequality follows directly from our argument above, since $v$ can be coupled with a vector $v_i$ whose coordinates are distributed according to Bernoulli$(p_i)$, so that we always have that $v$ dominates $v_i$ on every coordinate. The conclusion is that 
\[\Pb[\langle w, v\rangle\geq p]=\lim_{p_i\to p}\Pb[\langle w, v\rangle\geq p_i]\geq \lim_{p_i\to p}p_i=p\qedhere\]
\end{proof}

\begin{proposition}\label{prop:counterexample}
The answer to Question~\ref{q:prob} is negative for any $p\in (1/3, 1)$, with the exception of $p=1/2$. 
\end{proposition}
\begin{proof}
Let us first consider the case $1/3<p<1/2$. In this case, choosing $w=(1/3, 1/3, 1/3)$, we see that $\langle w, v\rangle\geq p$ if and only if $v$ has ones on at least $2$ coordinates. Thus,
\[\Pb[\langle w, v\rangle\geq p]=3p^2(1-p)+p^3=p(3p-2p^2).\]
It is not hard to verify that this quantity is smaller than $p$ for all $p\in (1/3, 1/2)$, since 
\[p-p(3p-2p^2)=p(1-3p+2p^2)=p(1-2p)(1-p)>0.\]
If $1/2<p<1$, we have a similar example, with $w=(1/2, 1/2)$. Again, we see that $\langle w, v\rangle\geq p$ if and only if $v$ has ones on at least $2$ coordinates. Thus,
\[\Pb[\langle w, v\rangle\geq p]=p^2<p.\qedhere\]
\end{proof}

\medskip
\textbf{Acknowledgements.} We would like to thank Byron Chin, Saba Lepsveridze and Alan Yan for their valuable comments on this paper.

\begin{thebibliography}{99}
\bibitem{BM87}
T.~Bier and N.~Manickam,
\emph{The first distribution invariant of the Johnson-scheme},
\emph{SEAMS Bull. Math.} \textbf{11} (1987), 61--68.

\bibitem{MM88}
N.~Manickam and D.~Mikl\'os,
\emph{On the number of nonnegative partial sums of a nonnegative sum},
In \emph{Colloq. Math. Soc. Janos Bolyai} \textbf{52} (1987), 385--392.

\bibitem{MS88}
N.~Manickam and N.~M.~Singhi,
\emph{First distribution invariants and EKR theorems},
\emph{J. Combin. Theory Ser. A} \textbf{48} (1988), 91--103.

\bibitem{AEFT13}
N.~Alon, Y.~Emek, M.~Feldman, and M.~Tennenholtz,
\emph{Adversarial leakage in games},
\emph{SIAM J. Discrete Math.} \textbf{27} (2013), 363--385.

\bibitem{AHS11}
N.~Alon, H.~Huang, and B.~Sudakov,
\emph{Nonnegative $k$-sums, fractional covers, and probability of small deviations},
\emph{J. Combin. Theory Ser. B}, \textbf{102} (2012), 784--796.

\bibitem{Tyom12}
M.~Tyomkyn,
\emph{An improved bound for the Manickam--Mikl\'os--Singhi conjecture},
\emph{European J. Combin.} \textbf{33} (2012), 27--32.

\bibitem{GM14}
S.~Greenberg and M.~Mohri,
\emph{Tight lower bound on the probability of a binomial exceeding its expectation},
\emph{Statistics \& Probability Letters} \textbf{86} (2014), 91--98.

\bibitem{PR16}
C.~Pelekis and J.~Ramon,
\emph{A lower bound on the probability that a binomial random variable is exceeding its mean},
\emph{Statistics \& Probability Letters} \textbf{119} (2016), 305--309.

\bibitem{Doe18}
B.~Doerr,
\emph{An elementary analysis of the probability that a binomial random variable exceeds its expectation},
\emph{Statistics \& Probability Letters} \textbf{139} (2018), 67--74.

\bibitem{Pok15}
A.~Pokrovskiy,
\emph{A linear bound on the Manickam--Mikl\'os--Singhi conjecture},
\emph{J. Combin. Theory Ser. A} \textbf{133} (2015), 280--306.

\end{thebibliography}


\end{document}
